\documentclass{2ssmeeting}

%%%%%%%%%%%%%%%%%%%%%%%%%%%%%%%%%%%%%%%%%%%%%%%%%%%%%%%%%%%%%%%%%%%%%%%%
%% POR FAVOR, NÃO FAÇA MUDANÇAS NESSE PADRÃO QUE ACARRETEM  EM
%% ALTERAÇÃO NA FORMATAÇÃO FINAL DO TEXTO
%%%%%%%%%%%%%%%%%%%%%%%%%%%%%%%%%%%%%%%%%%%%%%%%%%%%%%%%%%%%%%%%%%%%%%%%

%%%%%%%%%%%%%%%%%%%%%%%%%%%%%%%%%%%%%%%%%%%%%%%%%%%%%%%%%%%%%%%%%%%%%%%%
% POR FAVOR, ESCOLHA CONFORME O CASO
%%%%%%%%%%%%%%%%%%%%%%%%%%%%%%%%%%%%%%%%%%%%%%%%%%%%%%%%%%%%%%%%%%%%%%%%
\usepackage[brazil]{babel} % texto em português
% \usepackage[english]{babel} % texto em inglês

%\usepackage[latin1]{inputenc} % acentuação em Português ISO-8859-1
\usepackage[utf8]{inputenc} % acentuação em Português UTF-8
%%%%%%%%%%%%%%%%%%%%%%%%%%%%%%%%%%%%%%%%%%%%%%%%%%%%%%%%%%%%%%%%%%%%%%%%


%%%%%%%%%%%%%%%%%%%%%%%%%%%%%%%%%%%%%%%%%%%%%%%%%%%%%%%%%%%%%%%%%%%%%%%%
%% POR FAVOR, NÃO ALTERAR
%%%%%%%%%%%%%%%%%%%%%%%%%%%%%%%%%%%%%%%%%%%%%%%%%%%%%%%%%%%%%%%%%%%%%%%%
\usepackage[T1]{fontenc}
\usepackage{float}
\usepackage{graphics}
\usepackage{graphicx}
\usepackage{epsfig}
\usepackage{indentfirst}
\usepackage{amsmath, amsfonts, amssymb, amsthm}
\usepackage{url}
\usepackage{csquotes}
\setlength{\parindent}{0em}
% Ambientes pré-definidos

%%%%%%%%%%%%%%%%%%%%%%%%%%%%%%%%%%%%%%%%%%%%%%%%%%%%%%%%%%%%%%%%%%%%%%%%

\begin{document}

%%%%%%%%%%%%%%%%%%%%%%%%%%%%%%%%%%%%%%%%%%%%%%%%%%%%%%%%%%%%%%%%%%%%%%%%
% TÍTULO E AUTORES
%%%%%%%%%%%%%%%%%%%%%%%%%%%%%%%%%%%%%%%%%%%%%%%%%%%%%%%%%%%%%%%%%%%%%%%%
\title{Modelagem espaço-temporal bayesiana do risco de roubo de celulares nos municípios de São Paulo utilizando INLA: estimação e predição}

% Se necessário, você deve acrescentar ou remover autores no comando abaixo. Se os autores pertencerem a mesma instituição é possível colocar apenas uma vez a instituição.

\author{
    {\large Caio Gomes Alves}\thanks{c267583@dac.unicamp.br}\\
    {\small \textit{IMECC - UNICAMP}} \\
    {\large Pedro Francisco Godoy Bernardinelli}\thanks{p251570@dac.unicamp.br}\\
    {\small \textit{IMECC - UNICAMP}} \\
    {\large Rafael Ribeiro Santos}\thanks{r243464@dac.unicamp.br}\\
    {\small \textit{IMECC - UNICAMP}}
}
\criartitulo
%%%%%%%%%%%%%%%%%%%%%%%%%%%%%%%%%%%%%%%%%%%%%%%%%%%%%%%%%%%%%%%%%%%%%%%%

%%%%%%%%%%%%%%%%%%%%%%%%%%%%%%%%%%%%%%%%%%%%%%%%%%%%%%%%%%%%%%%%%%%%%%%%
% RESUMO
%%%%%%%%%%%%%%%%%%%%%%%%%%%%%%%%%%%%%%%%%%%%%%%%%%%%%%%%%%%%%%%%%%%%%%%%

\vspace{1cm}

\textbf{RESUMO}

\vspace{1cm}

Este trabalho apresenta uma modelagem espaço-temporal bayesiana para estimar o risco de roubo de celulares em 645 municípios do estado de São Paulo entre 2017 e 2024. Utilizando dados da Secretaria de Segurança Pública , a pesquisa empregou o método Integrated Nested Laplace Approximations (INLA) para ajustar um modelo de regressão Poisson, que inclui efeitos aleatórios espaciais e temporais, além de um termo de interação. O uso do INLA, que é mais rápido e preciso que o MCMC, permitiu obter as distribuições a posteriori marginais para cada parâmetro. A seleção do melhor modelo foi realizada com base no Deviance Information Criterion (DIC) e na minimização do erro quadrático médio de predição, resultando na escolha do modelo com interação do Tipo II, que aninha o efeito espacial no efeito temporal. Os resultados mostraram um declínio acentuado no risco de roubo de celulares durante os anos da pandemia de COVID-19, seguido por um crescimento posterior que se prevê continuar em 2025. Além disso, a análise espacial revelou um deslocamento do risco, que em 2017 estava mais concentrado no interior do estado e em 2024 se moveu para o litoral e a região metropolitana de São Paulo. A moda a posteriori foi usada como o valor predito para o risco, com mapas representando municípios mais seguros (risco relativo < 1) e mais perigosos (risco relativo > 1). O estudo conclui que a metodologia bayesiana com INLA é eficaz para analisar dados de criminalidade, fornecendo insights importantes para o combate ao crime.

\vspace{1.5em}

{\bf Palavras-chave:} Estatística Espacial; INLA; Modelos Mistos; Risco; Espaço-temporal.

\end{document}
